\chapter{\label{cha:discussion-conclusion}Discussion, conclusion, and future work}

\section{Summary of the thesis}
\label{sec:discussion-summary}

This thesis studied spectral bias in a controlled Neural Tangent Kernel (NTK)
setting with a one-hidden-layer ReLU network trained on functions defined with
inputs on the unit circle \(S^1\). The ideal reference point was the continuum
infinite-width limit with the uniform input distribution. In this limit, the
NTK is rotation-invariant, so the corresponding operator is diagonal in the
Fourier basis (the orthonormal basis of trigonometric functions \(\sin(k\theta),
\cos(k\theta)\)). This gives a clean description of learning via gradient
descent, i.e. each Fourier mode (specific harmonic i.e. a frequency \(k\)) in the target signal
decays at a rate determined by the corresponding kernel eigenvalue.

The central question was how far this description remains useful when the
idealized setting is made more realistic. The thesis considered two changes.
First, the continuum input distribution was replaced by finitely many sampled
points while keeping the infinite-width limit. Second, while keeping the
continuum input distribution, the infinite-width NTK was replaced by the NTK of
a finite-width network at initialization. Lastly, the finite-width NTK was
allowed to evolve during training.

The answer is not the same in all three cases. With finite data, the exact
diagonalization of the NTK in the Fourier basis is lost, but the learning
dynamics remain close to the continuum predictions on low-frequencies when the
sample size is large enough. More precisely, for the low-frequency subspace
\(\mathcal H_K\) composed of \(0,1, \ldots K\) frequencies, the quantities that
measure how error between the sampled operator and the diagonal operator all
become small, i.e. the Fourier basis functions, when evaluated on the sampled inputs, become
approximately orthonormal under the empirical inner product, the
sampled operator restricted to \(\mathcal H_K\) becomes closer to the predicted
diagonal operator, and applying the sampled operator to a retained Fourier mode
becomes closer to multiplying that mode by its predicted eigenvalue. For fixed
\(K\), these deviations decrease at \(O(n^{-1/2})\), up to logarithmic
factors, and the probability that they exceed any fixed threshold decays
exponentially in \(n\).

With a finite-width neural network of width \(m\), the NTK at initialization is a random approximation
of the infinite-width operator. On the same fixed low-frequency subspace, the
difference between the finite-width operator and the infinite-width diagonal
Fourier operator decreases at the \(m^{-1/2}\) scale, again up to logarithmic
factors. Equivalently, the probability that they exceed and fixed threshold decays
exponentially in the width \(m\).

When the finite-width NTK evolves during training, the Fourier modes are no
longer fixed eigenvectors of the operator. The kernel can both rotate
eigenspaces away from the Fourier subspaces and change how much strength it
assigns to different frequencies. The Fourier basis is therefore no longer an
exact diagonalization of the dynamics, but it remains useful for measuring how
training changes residual decay, subspace alignment, and kernel strength across
frequencies.

The main conclusion is that Fourier modes remain useful for understanding
spectral bias beyond the ideal continuum model, but in different ways in
different regimes. In the continuum infinite-width model, the Fourier modes are
exact eigenvectors of the NTK operator, so each frequency has an exact predicted
learning rate. Under finite sampling and finite width, this picture
remains approximately accurate on fixed low-frequency subspaces. During
finite-width training, the NTK can evolve and move away from this fixed
Fourier-diagonal structure. In that regime, the Fourier subspaces are not exact
eigenspaces, but they still provide reference directions for tracking how
training changes residual decay, subspace alignment, and kernel strength.

\section{How finite sampling changes the Fourier-mode picture}
\label{sec:discussion-finite-data}

The finite-sampling part of the thesis makes precise what changes when the
uniform distribution on \(S^1\) is replaced by finitely many uniformly sampled
points. In the continuum setting, each Fourier mode evolves independently under
the NTK operator. After sampling, the retained Fourier modes are evaluated only
at the sampled points, and the sampled operator is a finite matrix rather than
the continuum integral operator. The question is whether the sampled dynamics
still stays close to the diagonal Fourier prediction on the low-frequency modes
that we retain.

\subsection{What the finite-sample bounds control}
\label{subsec:discussion-finite-data-structure}

In the notation described in the chapter \ref{cha:methodology}, the finite-sample analysis controls
three quantities:
\[
	G-I_d,
	\qquad
	H-\Lambda^{(n)},
	\qquad
	A\Phi-\Phi\Lambda^{(n)}.
\]
These quantities measure three concrete ways in which sampling can break the
continuum Fourier picture.

First, \(G-I_d\) measures how far the retained Fourier modes are from being
orthonormal after they are evaluated on the sampled points. Second,
\(H-\Lambda^{(n)}\) measures how far the sampled operator, restricted to the
retained low-frequency block, is from the predicted diagonal operator. Here
\(\Lambda^{(n)}\) contains the finite-sample corrected eigenvalues. Thus this
quantity directly measures whether the sampled operator still acts mainly by
rescaling each retained Fourier mode. Third, \(A\Phi-\Phi\Lambda^{(n)}\)
measures the error made when applying the sampled operator to the retained
Fourier modes. If this term were zero, then applying the sampled operator to
each retained Fourier mode would be exactly the same as multiplying that mode
by its predicted finite-sample eigenvalue.

Theorem~\ref{thm:results-finite-sample-concentration} and
Figure~\ref{fig:results-finite-sample-diagnostics} show that all three errors
decrease as the sample size \(n\) increases, for fixed low-frequency cutoff
\(K\). For fixed \(K\), the typical size of these deviations decreases at the
\(n^{-1/2}\) scale, up to logarithmic factors. Equivalently, for any fixed error
threshold, the probability of exceeding that threshold decays exponentially in
\(n\).

This is the precise sense in which the sampled problem remains close to the
continuum Fourier prediction on a fixed low-frequency subspace. The result does
not say that the sampled operator is exactly diagonal but that the
deviation from the diagonal prediction becomes small with high probability as
the number of sampled points increases.

This connects to the broader operator-spectrum view of spectral bias. In the
uniform population setting, the Fourier modes diagonalize the kernel operator.
For non-uniform input distributions, the relevant eigenfunctions depend on the
input density rather than being the standard global Fourier modes
\citep{basri2020frequencybiasneuralnetworks,bietti2019inductivebiasneuraltangent,
	cao2020understandingspectralbiasdeep}. The finite-sample result in this thesis
sits between these two settings: the samples are drawn from the uniform
distribution, but the empirical operator is still a finite random approximation
of the uniform population operator.

\subsection{Consequences for residual dynamics}
\label{subsec:discussion-finite-data-dynamics}

The residual plots show this effect directly. During the first part of training,
the observed Fourier-mode amplitudes closely follow the decay rate as predicted by the
eigenvalue. This means that the main ordering of learning speeds is
still the one predicted by the diagonal Fourier model: modes with larger
predicted eigenvalues decay faster, and modes with smaller predicted eigenvalues
decay more slowly.

The reason the curves eventually deviate is visible from the operator
decomposition
\[
	A\Phi
	=
	\Phi\Lambda^{(n)}
	+
	\bigl(A\Phi-\Phi\Lambda^{(n)}\bigr).
\]
The first term is the diagonal prediction. If this were the whole operator
action, then each retained Fourier mode would evolve independently at its
predicted rate. The second term is the part of the sampled operator action that
is not captured by this diagonal prediction.

To see the effect on the dynamics, suppose for the moment that the sampled
residual lies mainly in the retained Fourier span, so that
\[
	r_t \approx \Phi \alpha_t .
\]
Then applying the sampled operator gives
\[
	A r_t
	\approx
	A\Phi\alpha_t
	=
	\Phi\Lambda^{(n)}\alpha_t
	+
	\bigl(A\Phi-\Phi\Lambda^{(n)}\bigr)\alpha_t .
\]
The first term gives the ideal mode-wise decay. The second term is an additional
forcing term coming from the finite-sample error. It can move mass between the
retained Fourier coordinates and can also push the residual away from the
diagonal prediction.

Early in training, the residual components are still large, so the diagonal
term is the dominant part of the motion. This is why the observed curves follow
the predicted rates well at the beginning. Later, after some modes have already
decayed by several orders of magnitude, the remaining error term can become
visible in the normalized curves. A small absolute error can look large once the
quantity being plotted is already very small. This explains why the diagonal
prediction captures the initial decay rates better than the late-time behavior.

This also explains the difference between the \(n=4096\) and \(n=32768\)
experiments. With more samples, the non-diagonal part of the sampled operator is
smaller, so the observed mode curves stay close to the diagonal prediction for
longer.

\subsection{What the preconditioning experiment shows}
\label{subsec:discussion-finite-data-preconditioning}

The preconditioning experiment gives another way to test the same point. The
baseline sampled dynamics show different decay speeds across Fourier modes. The
question is how much of this separation is explained simply by the predicted
Fourier eigenvalues.

The theory preconditioner uses only the diagonal prediction
\(\Lambda^{(n)}\). It rescales the retained Fourier coefficients by the inverse
predicted eigenvalues, but it does not use the observed sampled block
\(C_n^{(K)}=G^{-1}H\). The empirical preconditioner does use
\(C_n^{(K)}\), so it corrects both the diagonal scaling and the finite-sample
mixing inside the retained block.

Figures~\ref{fig:results-preconditioned-mode-decay},
\ref{fig:results-theory-precond-spectrum}, and
\ref{fig:results-preconditioner-gap} show that the theory preconditioner already
removes most of the rate separation on the retained block. This means that most
of the observed separation is explained by the predicted Fourier eigenvalues.
The empirical preconditioner gives the flattest behavior because it also uses
the observed sampled block and therefore corrects the remaining mixing caused by
the particular sampled points.

This should be interpreted as a diagnostic, not as a proposed training method.
It shows that, in this experiment, the diagonal Fourier prediction explains most
of the finite-sample learning-rate separation on the retained modes. The gap
between the theory and empirical preconditioners measures what is left after
that diagonal scaling has been accounted for.

\subsection{Takeaway for finite sampling}
\label{subsec:discussion-finite-data-summary}

Finite sampling destroys exact Fourier diagonalization, but it does not destroy
the usefulness of the diagonal Fourier prediction on fixed low-frequency
blocks. As \(n\) increases, the retained Fourier modes become closer to
orthonormal on the sampled points, the sampled operator becomes closer to the
predicted diagonal operator on the retained block, and applying the sampled
operator to retained Fourier modes becomes closer to multiplying each mode by
its predicted eigenvalue.

This is why the corrected eigenvalues predict the main early-time residual
decay. The later deviations are caused by the remaining part of the sampled
operator that is not captured by the diagonal Fourier prediction.

\section{Role of the trainable bias}
\label{sec:discussion-bias}

The trainable hidden bias changes which Fourier modes are present in the
continuum fixed-kernel dynamics. Corollary~\ref{cor:results-bias-spectrum}
shows that, without the bias contribution, the odd modes \(k\ge 3\) have zero
eigenvalue. With the bias contribution, these modes have nonzero but small
eigenvalues. Thus the bias determines whether these modes can be learned at all
in the fixed-kernel continuum dynamics.

This is consistent with earlier harmonic analyses of two-layer ReLU networks on
the sphere and circle, where bias-free two-layer models do not express odd
harmonics beyond the first frequency
\citep{basri2019convergence,basri2020frequencybiasneuralnetworks}. In the
present calculation, the same effect appears directly in the NTK decomposition:
removing the trainable hidden-bias contribution sets the odd eigenvalues
\(k\ge 3\) to zero, while adding it makes them nonzero. Thus
Figure~\ref{fig:results-continuum-odd-modes} distinguishes modes that are
present with small eigenvalues from modes that are absent because the
corresponding eigenvalues vanish.

This point is useful because it shows that the exact architecture matters even
inside the simplified NTK setting. The continuum Fourier picture is not just a
generic statement that low frequencies are learned faster. It also depends on
which parameter directions are included in the tangent kernel. Training the
hidden biases adds directions to the tangent feature space, and those directions
make additional odd Fourier modes visible to the fixed-kernel dynamics.

\section{How finite width changes the Fourier description}
\label{sec:discussion-finite-width}

The finite-width part of the thesis keeps the continuum input distribution on
\(S^1\) fixed and asks what changes when the infinite-width tangent kernel is
replaced by the tangent kernel of a finite-width network. There are two
separate issues. The first is the frozen tangent kernel at initialization. This
kernel is random because it is built from finitely many neurons. The second is
the evolving tangent kernel during training. This kernel can move away from its
initial value, so it need not remain close to the fixed Fourier-diagonal
reference.

\subsection{Frozen finite-width kernels}
\label{subsec:discussion-finite-width-frozen}

The finite-width theorem concerns the frozen tangent operator at initialization.
Theorem~\ref{thm:results-finite-width-frozen-concentration} shows that
\[
	\mathbb E[B_m^{(K)}]=\Lambda_K
\]
and that \(B_m^{(K)}\) concentrates around \(\Lambda_K\) as \(m\) increases.
Thus, on a fixed low-frequency block, finite width acts as a random perturbation
of the continuum Fourier operator.

Figure~\ref{fig:results-finite-width-concentration} supports this conclusion.
The median and \(95\%\) quantile of
\(\|B_m^{(K)}-\Lambda_K\|_{\max}\) both decrease as \(m\) increases. The
theoretical envelope is conservative, but the qualitative trend matches the
concentration statement. This agrees with the role of the infinite-width NTK as
a useful baseline, while also showing why finite-width fluctuations can matter
quantitatively
\citep{jacot2018ntk,lee2020wide,bordelon2023dynamicsfinitewidthkernel,vyas2022limitationsntkunderstandinggeneralization}.

\subsection{Eigenspace alignment at finite width}
\label{subsec:discussion-finite-width-alignment}

The subspace-alignment plots give an eigenspace-level view of the same
finite-width effect. Figures~\ref{fig:results-finite-width-projector-error} and
\ref{fig:results-finite-width-principal-angles} show that the matched empirical
eigenspaces become more aligned with the corresponding Fourier subspaces as
\(m\) increases. This matches the trend predicted by the concentration result for
\(B_m^{(K)}-\Lambda_K\).

The ordering across frequencies is not perfectly monotone. In particular, the
\(k=2\) Fourier plane is better aligned than some lower-frequency blocks in the
projector-error plot. One possible reason is that \(k=0\) and \(k=1\) are
structurally special in this model. The constant mode is one-dimensional and
receives a direct contribution from the trainable-bias part of the NTK. The
first frequency is tied to the coordinate functions in the embedding
\(x(\varphi)=(\cos\varphi,\sin\varphi)\), and it is the only odd frequency
already present in the no-bias kernel:
\[
	\lambda_1^{(\mathrm{nb})}=\frac14,
	\qquad
	\lambda_k^{(\mathrm{nb})}=0
	\quad \text{for odd } k\ge 3.
\]
Thus \(k=0\) and \(k=1\) are not generic representatives of the
higher-frequency pattern.

The main conclusion from the alignment diagnostics is the width trend. For each
displayed frequency, the matched empirical subspace becomes more aligned with
the corresponding Fourier subspace as \(m\) increases.

\subsection{Evolving kernels during training}
\label{subsec:discussion-finite-width-evolving}

The evolving-kernel experiments add an empirical part to the finite-width
answer. The loss curves show a width-dependent crossover: the evolving tangent
kernel improves the final objective at widths \(128\) and \(256\), gives a
smaller improvement at width \(512\), and loses its advantage at width \(1024\).

The joint diagnostics suggest that this behavior is not explained by improved
Fourier geometry. The subspace cosine similarity plots show that several
higher-frequency subspaces become less aligned with their Fourier reference
subspaces during evolving training. In this sense, kernel evolution does not
make the tangent kernel more Fourier-diagonal.

The spectral-strength plots point to a different effect. During evolving
training, the average kernel strength increases in the lowest Fourier blocks,
especially \(k=0\) and \(k=1\), with a smaller increase for \(k=2\). This
matches the mode-wise residual plots: at small width, the modes that improve
most under evolving dynamics are the low-frequency target modes. Thus the
benefit of kernel evolution in these experiments appears to come from increasing
the strength of the operator on the low-frequency blocks, rather than from
preserving the Fourier eigenspaces more accurately.

The drift diagnostic
\[
	\|C_t^{(K)}-C_0^{(K)}\|_F
\]
supports the same width-dependent picture. It measures how much the projected
low-frequency tangent-kernel block moves away from initialization. The drift is
largest at small width and much smaller at large width. At small width, this
larger change is accompanied by increased low-frequency strength and improved
loss decay. At large width, the frozen kernel is already closer to the
continuum Fourier baseline, and the extra change from kernel evolution is
smaller. This helps explain why the evolving-kernel advantage shrinks with
width.

The amplitude-variant experiments further show that the low-frequency strength
increase is not specific to the original target amplitudes. Even when the
non-constant modes have equal amplitudes, or when the larger amplitudes are
assigned to higher frequencies, the evolving kernel still increases cumulative
low-frequency strength. This does not prove a general causal mechanism, but it
does make the spectral-strength interpretation more robust than an explanation
based only on the original target.

\subsection{Summary for finite width}
\label{subsec:discussion-finite-width-summary}

For the frozen tangent operator, finite width preserves the continuum Fourier
block up to a random perturbation that decreases with \(m\). The concentration
theorem proves this at the matrix level, and the eigenspace diagnostics show the
same trend geometrically: as width increases, the matched empirical subspaces
become closer to the corresponding Fourier frequency subspaces.

For the evolving tangent kernel, the thesis gives an empirical answer rather
than a theorem. Kernel evolution changes the finite-width dynamics in a
width-dependent way. At small width, it improves the final loss and accelerates
the decay of several low-frequency target modes. The diagnostics suggest that
this improvement is associated with increased low-frequency spectral strength,
not with improved Fourier-subspace alignment. At larger width, the frozen kernel
is already a stronger approximation to the continuum Fourier operator, and the
advantage of the evolving kernel decreases or disappears.

Thus the answer for finite width is mixed. In the frozen setting, the Fourier
description is recovered as width increases. In the evolving setting, the
Fourier subspaces remain useful diagnostics, but the kernel is no longer simply
a fixed Fourier-diagonal object. It changes during training by both
redistributing strength across frequencies and moving some eigenspaces away from
the Fourier reference. A comparable non-asymptotic theory for this evolving
regime remains open.

\section{Why high frequencies are more sensitive}
\label{sec:discussion-high-frequency-sensitivity}

One theme appearing in both the finite-sample and finite-width results is that
higher frequencies are more fragile. In the continuum model, high-frequency
modes are learned slowly because their eigenvalues are small. The same small
eigenvalues also make these modes more sensitive to finite-sample and
finite-width errors: an operator error that is small in absolute size can still
be large compared with the learning rate of a high-frequency mode.

From Theorem~\ref{thm:results-continuum-spectrum},
\[
	\lambda_k = O(k^{-2})
	\qquad (k\to\infty),
\]
for both even and odd frequencies. Therefore, as \(k\) increases, the same
absolute operator perturbation becomes larger relative to the continuum decay
rate of that frequency.

The same issue appears at the eigenspace level. For each \(k\ge 1\), the cosine
and sine modes form a two-dimensional Fourier plane. At finite width, the
matched empirical eigenspace can rotate away from this plane. This rotation is
more visible when neighbouring Fourier eigenvalues are close together, which is
more common at higher frequencies. In the evolving finite-width experiments,
higher-frequency subspaces also show larger losses of Fourier alignment during
training.

Thus, in this model, high frequencies are not only slower to learn; their
Fourier structure is also less stable under sampling, finite-width fluctuations,
and kernel evolution. This is consistent with empirical observations that
high-frequency components are learned later and are more sensitive to the data
distribution and to perturbations
\citep{rahaman2019spectralbiasneuralnetworks,basri2019convergence,basri2020frequencybiasneuralnetworks}.

\section{Limitations}
\label{sec:discussion-limitations}

\subsection{Fixed low-frequency blocks}
\label{subsec:discussion-limitations-fixed-block}

The concentration results are fixed-block results. They control a prescribed
low-frequency Fourier subspace with fixed cutoff \(K\). They do not give a
uniform statement over all Fourier modes, and the bounds are not designed to
remain sharp as \(K\) grows with \(n\) or \(m\). This is important because the
high-frequency eigenvalues and eigengaps become small, so controlling higher
and higher frequencies would require increasingly accurate operator
approximations.

\subsection{Conservative concentration bounds}
\label{subsec:discussion-limitations-concentration}

The numerical constants in the concentration bounds are conservative. In the
finite-sample theorem, the proof controls entries of the relevant matrices and
then applies union bounds over the retained block. The bound for
\(H-\Lambda^{(n)}\) uses a bounded-differences argument, and the bound for
\(A\Phi-\Phi\Lambda^{(n)}\) uses a conditional Bernstein estimate followed by a
union bound over sample-mode pairs. In the finite-width theorem, each projected
entry is controlled using a coarse tail envelope before another union bound is
applied. These steps are sufficient to prove concentration, but they discard
matrix structure and variance information. This explains why the theoretical
envelopes in Figures~\ref{fig:results-finite-sample-diagnostics} and
\ref{fig:results-finite-width-concentration} are much larger than the observed
errors.

In the finite-width concentration plot, the theorem contains an unspecified
universal constant \(c>0\). This constant comes from Bernstein's inequality for
averages of centered sub-exponential random variables
\citep{Vershynin_2018a}. In explicit proofs of such inequalities, constants of
order \(1/4\) often appear after converting a \(\psi_1\)-norm bound into a
moment generating function bound. The plotted choice \(c=0.2\) is a conservative
representative value, not an optimized constant for this problem. Changing this
constant shifts the envelope but does not change the predicted dependence on
\(m\).

A sharper analysis would need to use the structure of the objects being
controlled. For finite sampling, \(H\) is a bounded second-order statistic, so
concentration tools for U- or V-statistics may give tighter estimates. For both
finite sampling and finite width, direct matrix concentration in operator norm
would be preferable to entrywise control followed by norm conversion.

\subsection{Limits of the dynamics claims}
\label{subsec:discussion-limitations-dynamics}

The finite-sample dynamics experiments support the diagonal prediction, but the
theory does not give a complete time-uniform residual bound. The experiments
show that the predicted eigenvalues explain the dominant early-time behavior.
At later times, some residual components have already become very small, so even
a small difference between the sampled operator and the diagonal prediction can
become visible in relative terms. The current bounds identify the source of this
effect, but they do not give a sharp description of the full residual trajectory
for all times.

The preconditioning experiment is also a diagnostic rather than a proposed
training method. The theory preconditioner tests whether the observed rate
separation is mainly explained by the diagonal eigenvalue scaling. The empirical
preconditioner tests how much is gained by using the actual sampled
low-frequency block. These experiments support the interpretation of the
sampled dynamics, but they do not establish a general preconditioning algorithm
or a generalization guarantee.

\subsection{Limits of the finite-width alignment diagnostics}
\label{subsec:discussion-limitations-alignment}

The projector-error and principal-angle plots are empirical diagnostics of the
frozen finite-width operator. They compare Fourier frequency subspaces with
matched empirical eigenspaces of a discretized operator. This is useful for
checking whether the eigenspaces become more Fourier-like as width increases,
but it is not a separate theorem.

The matching procedure also introduces a practical limitation. For each
frequency \(k\), the empirical eigenspace is chosen by overlap with the
corresponding Fourier subspace. This makes the comparison interpretable, but it
also means that the diagnostic depends on the selected discretization grid and
on the eigenspace-matching rule. The most useful conclusion from these plots is
the trend of improved alignment as width increases.

\subsection{Limits of the evolving-kernel experiments}
\label{subsec:discussion-limitations-evolving}

The finite-width theorem concerns only the tangent operator at initialization.
It does not prove a bound for the evolving tangent kernel during training. The
frozen result therefore controls one source of finite-width error, namely random
initialization of the tangent kernel, but not the time-dependent change of the
kernel along the training trajectory.

The evolving-kernel section provides an empirical comparison between frozen and
evolving dynamics. It tracks the loss, mode-wise residual decay,
Fourier-subspace alignment, spectral strength, and block drift. These
diagnostics show a coherent width-dependent pattern, but they do not amount to a
time-dependent concentration theorem. A full theory would need to control the
evolving block \(C_t^{(K)}\), its drift from \(C_0^{(K)}\), and its action on the
current residual.

There is also no joint theorem combining finite sampling, finite width, and
training time. The finite-sample analysis studies the infinite-width sampled
operator. The finite-width theorem studies the continuum frozen operator. In
actual finite training runs, sampling error, finite-width initialization error,
and kernel evolution occur together.

\subsection{Model assumptions}
\label{subsec:discussion-limitations-model}

The analysis relies on the special structure of \(S^1\) with the uniform
measure. In this setting, the limiting kernel is rotation-invariant, so the
corresponding operator is a convolution operator and the Fourier basis
diagonalizes it exactly. Related harmonic descriptions of neural tangent kernels
on the circle or sphere appear in
\citet{basri2019convergence,basri2020frequencybiasneuralnetworks,bietti2019inductivebiasneuraltangent,cao2020understandingspectralbiasdeep}.
This structure is no longer automatic once the input distribution is changed.
For example, under a non-uniform density on the circle, the integral operator is
weighted by the density and is no longer diagonal in the global Fourier basis;
the relevant eigenfunctions become distribution-dependent, as studied by
\citet{basri2020frequencybiasneuralnetworks}. If the input domain is changed,
the appropriate basis also changes and may not be available in closed form.

The architecture and scaling assumptions are also part of the result. The thesis
uses a one-hidden-layer ReLU network in NTK scaling. The ReLU activation is
important because its Gaussian averages have explicit arc-cosine formulas
\citep{cho2009kernel}, which makes the kernel and its Fourier coefficients
computable in closed form. The \(1/\sqrt m\) scaling is important because it
places the model in the NTK or lazy-training regime, where the infinite-width
limit is described by a fixed kernel
\citep{jacot2018ntk,lee2020wide,golikov2022neuraltangentkernelsurvey}.
Changing the activation, adding or removing bias terms, changing the
parameterization, or moving to a feature-learning scaling changes the limiting
kernel and therefore changes the spectrum that drives the dynamics. Such
feature-learning regimes are better described by mean-field or related
infinite-width limits
\citep{Mei2018meanfieldview,chizat2018globalconvergencegradientdescent,sirignano2019meanfieldanalysisneural,Rotskoff_2022}.
Thus the Fourier conclusions should be understood as conclusions about this
explicit NTK model, not as architecture-independent statements about neural
networks in general.

\section{Future work}
\label{sec:future-work}

The most direct next step is to develop theory for the evolving finite-width
tangent kernel. The experiments in this thesis compare frozen and evolving
dynamics, but the analysis remains empirical. A natural goal is to control the
time-dependent tangent-kernel block, its drift from initialization, and its
effect on the residual dynamics.

A second direction is to prove a joint result that includes finite sampling,
finite width, and training time. The present thesis studies finite sampling and
frozen finite width separately, and studies the evolving finite-width kernel
empirically. In actual training runs, all three effects occur at once.

A third direction is to sharpen the concentration bounds. The current proofs use
entrywise concentration, coarse tail bounds, and union bounds. These arguments
are sufficient to prove the correct decreasing trends, but they are conservative
as numerical estimates. Direct operator-norm bounds, variance-sensitive
estimates, or concentration inequalities for U- and V-statistics may give bounds
closer to the observed errors.

A fourth direction is to study the alignment diagnostics more directly. The
projector-error and principal-angle plots suggest that different Fourier
subspaces can have different finite-width stability. A more detailed analysis
could measure which Fourier blocks are most strongly mixed by the finite-width
operator and how this depends on the eigenvalues and gaps of the continuum
operator.

Finally, the analysis should be extended beyond the uniform circle setting.
Non-uniform input densities, higher-dimensional spheres, other activations, and
different parameterizations all change the relevant kernel operator and its
spectrum. Studying these cases would clarify which conclusions are special to
the explicit Fourier model on \(S^1\), and which reflect a broader mechanism
behind spectral bias.

\section{Conclusion}
\label{sec:conclusion}

The main conclusion of the thesis is that the Fourier description of spectral
bias remains a useful organizing principle beyond the ideal continuum
infinite-width model. It is exact in the continuum setting, approximately
accurate for finite data on fixed low-frequency blocks, and accurate for the
frozen finite-width tangent kernel as width increases. For the evolving
finite-width kernel, it is no longer a fixed diagonal description, but it remains
useful for tracking how training changes residual decay, Fourier alignment, and
kernel strength across frequencies.

The evolving finite-width kernel is different. During training, the tangent
kernel can change its alignment with the Fourier subspaces and redistribute its
strength across frequencies. In the experiments studied here, this evolution
improves optimization at small width mainly by increasing strength in the
low-frequency blocks, even though some higher-frequency subspaces become less
Fourier-aligned. At larger width, the frozen kernel is already closer to the
continuum reference, and the advantage of kernel evolution decreases.

Thus the thesis gives a controlled picture for finite data and frozen finite
width, and an empirical picture for evolving finite width. Developing a
comparable theory for the evolving tangent kernel remains the main open problem.
